\documentclass[10pt,a4paper]{amsart}
\usepackage[margin=19mm]{geometry}
\usepackage{amsmath,amssymb,mathtools}
\usepackage[scr=rsfs,cal=euler]{mathalfa}
\usepackage{xfrac}
\usepackage[unicode,hidelinks,bookmarks=false]{hyperref}
\newcommand{\ie}{i.e.\ }
\newcommand{\cf}{cf.\ }
\newcommand{\resp}{respectively\ }
\newcommand{\Subsection}{\subsection}
\providecommand{\nobreakdash}{\nobreak\hbox{-}}
\setlength{\emergencystretch}{2em}
\multlinegap0pt
\usepackage{xcolor}
\usepackage{amssymb}
\usepackage[normalem]{ulem}
\usepackage{enumerate}
\newtheorem{lem}{Lemma}
\newtheorem{thm}{Theorem}
\def\F{{\mathbb F}}
\title{Enumeration of certain permutation group polynomials}
\author{Sartaj Ul Hasan, Ramandeep Kaur, Hridesh Kumar}
\date{Source: arXiv:2608.28118v1 (CC BY 4.0)}
\begin{document}
\maketitle

\noindent\emph{Source and licence.} Enumeration of certain permutation group polynomials. Version arXiv:2608.28118v1 (CC BY 4.0), distributed by arXiv under the Creative Commons Attribution 4.0 International licence, http://creativecommons.org/licenses/by/4.0/, as recorded in the arXiv licence metadata for this exact version and shown on the abstract page https://arxiv.org/abs/2608.28118v1. Full licence notice: \texttt{licenses/arxiv-2608.28118-CC-BY-4.0.txt}.\par\smallskip

\noindent\textit{Editorial scope.} Counted unit: the article's lem Atleast\_lemma\_T31 (source0.tex lines 844--853). The article's complete original proof of it is delivered with it (source0.tex lines 854--1113, one \texttt{proof} environment of 260 lines). No mathematics is written, repaired or re-derived: the statement and the proof are the article's own lines, in the article's own notation, and no retained line is substituted or re-flowed.  Retained uncounted as the source-local context the proof needs: lines 240--256; lines 259--266; lines 268--275; lines 384--386; lines 495--497; lines 818--820.  Nothing was omitted from any retained span, so the package carries no omission record.  No retained line contains a citation, so the wrapper carries no bibliography and no bibliography entry.  No retained line contains a figure environment, an image-inclusion command or drawing code, so no image is delivered and the package declares no asset dependency.  Editorial formatting only, in addition: an AMS \texttt{amsart} preamble that restates the article's own declarations and macros used by the retained lines, namely the environments \texttt{lem}, \texttt{thm} on separate counters, together with the standard packages those lines need.  The retained environments print in this document as Lemma 1, Theorem 1 in order of appearance, whereas the article prints the same items with the numbering of its own full document.  The complete native source of the article as distributed in the arXiv e-print for this exact version is bundled unchanged as \texttt{sources/208762/source0.tex}.
 \begin{lem}\label{RR1}
 Let $q=p^3$, where $p$ is an odd prime. Let $\F_q = \{ c_0, c_1, \ldots, c_{q-1}\}$ be the finite field with $q$ elements,
\[
a=(c_{0},c_{1},\ldots,c_{p^2-1})(c_{p^2},c_{p^2+1},\ldots,c_{2p^2-1})\cdots (c_{p^2(p-1)},c_{p^2(p-1)+1},\ldots,c_{p^3-1}), 
\]
 and 
\[
b=b_{0}b_{1}^{1+p}b_2^{1+2p}\cdots b_{p-1}^{1+(p-1)p}
\]
be permutations of $\F_q$, where 
\begin{equation*}
\begin{split}
b_{k}=(c_{k},c_{p^2+k},\ldots,c_{(p-1)p^2+k},c_{p+k},c_{p+p^2+k},\ldots,c_{p+(p-1)p^2+k},\ldots,c_{(p-1)p+k},c_{(p-1)p+p^2+k},\ldots,& \\ c_{(p-1)p+(p-1)p^2+k}). 
\end{split}
\end{equation*}
Then $ab=b^{p^2-p+1}a$.
 \end{lem}

 \begin{equation}\label{e21}
a(c_{t})=\begin{cases}

       c_{t+1} & \text{ if } t < r_tp^{2}-1, \\
       c_{t+1-p^2} & \text{ if } t=r_tp^{2}-1
       	\end{cases}
=c_{(r_t-1)p^{2}+(t+1)\pmod {p^{2}}}.
\end{equation}

\begin{equation}\label{e22}
b_k(c_{t})=\begin{cases}
       c_{t} & \text{ if } t \not \equiv k  \pmod  p, \\
       c_{t+p^2} & \text{ if } t < p^{3}-p^2 \text{ and } t \equiv k  \pmod  p, \\
       c_{t-p^3+p^2+p} & \text{ if } p^{3}-p^2 \leq  t < p^{3}-p \text{ and } t \equiv k  \pmod  p,\\
       c_{t-p^3+p} & \text{ if }    t \geq p^{3}-p \text{ and } t \equiv k \pmod  p.
       	\end{cases}
\end{equation}

\begin{equation}\label{tempe4}
a^{d}(c_t)=c_{(r_t-1)p^2+(t+d) \pmod {p^2}}.
\end{equation}

\begin{thm}\label{T31}
Let $q=p^3$, where $p$ is an odd prime, and let $\F_q=\{c_0,c_1,\ldots,c_{q-1}\}$ denote the finite field with $q$ elements. Suppose that $a$ and $b$ defined as in Lemma~\ref{RR1} are permutations of $\F_q$. Then $G=\langle a,b\rangle$ is a subgroup of $\mathfrak{S}_q$ of order $q$, and every non-identity element of $G$ has no fixed point in $\F_q$.
\end{thm}

\begin{lem}\label{Atmost_lemma_T31}
Let $q=p^3$, where $p$ is an odd prime, and $G=\langle a,b \rangle $ be the subgroup of $\mathfrak{S}_q$ defined as in Theorem \ref{T31}. Moreover, let $i_1,i_2 \in  \{0,1,\ldots,p-1\}$ and $j_1,j_2 \in \{0,1,\ldots,p^{2}-1\}$ be fixed integers. Then $|\mathcal{N}(i_1,j_1,i_2,j_2)| \leq  p^3$.
\end{lem}

\begin{lem}\label{Atleast_lemma_T31}
Let $q=p^3$, where $p$ is an odd prime, and $G=\langle a,b \rangle$ be a subgroup of $\mathfrak{S}_q$ defined as in Theorem \ref{T31}. Furthermore,  let $i_1,i_2 \in \{0,1,\ldots,p-1\}$ and $j_1,j_2 \in \{0,1,\ldots,p^{2}-1\}$. Then $\mathcal{N}(i_1,j_1,i_2,j_2) \neq \emptyset$ if and only if 
\begin{equation*}
\begin{split}
& j_1+i_1 \not \equiv  0  \pmod  p,\\& 
  j_2+i_2 \not \equiv  0 \pmod  p,
\end{split}
\end{equation*}
 $i_2\equiv i_1-1 \pmod  p$ and $j_2 \equiv j_1+1 \pmod  p$. In this case,  $|\mathcal{N}(i_1,j_1,i_2,j_2)|=p^3$.
\end{lem}

\begin{proof}
We first assume that $\mathcal{N}(i_1,j_1,i_2,j_2) \neq \emptyset$, where $i_1,i_2 \in \{0,1,\ldots,p-1\}$ and $j_1,j_2 \in \{0,1,\ldots,p^{2}-1\}$. We will show that $ j_1+i_1 \not \equiv  0  \pmod  p$, $j_2+i_2 \not \equiv  0 \pmod  p$, $i_2\equiv i_1-1 \pmod  p$ and $j_2 \equiv j_1+1 \pmod  p$. Let $h \in \mathcal{N}(i_1,j_1,i_2,j_2)$. Thus, we have 
\begin{equation}\label{s41}
\begin{split}
&hah^{-1}=b^{j_1}a^{i_1}\\
&hbh^{-1}=b^{j_2}a^{i_2},
\end{split}
\end{equation}
which gives $\mid b^{j_1}a^{i_1}\mid =\mid b^{j_2}a^{i_2}\mid =p^2$ as $|a|=|b|=p^2$. Therefore, we have to exclude all the possibilities of $i_1,j_1,i_2,j_2$ such that $(b^{j_1}a^{i_1})^p=I$ or $(b^{j_2}a^{i_2})^p=I$. We make use of the relation $ab=b^{p^2-p+1}a$ and do some simplifications to obtain 
\[
(b^{j_1}a^{i_1})^p=b^{j_1N+i_1p},
\]
where 
$$
N=\begin{cases} \dfrac{(p^2-p+1)^{i_1p}-1}{(p^2-p+1)^{i_1}-1} & \text{ if } i_1 \neq 0,\\p 
 & \text{ if } i_1=0.\end{cases}$$
This implies that  $(b^{j_1}a^{i_1})^p=I$ if and only if $j_1N+i_1p \equiv 0 \pmod {p^2}$.
 It is trivial that $N \equiv p \pmod {p^2}$ if $i_1=0$. Now we shall show that $N \equiv p \pmod {p^2}$ also holds for $1 \leq i_1 \leq p-1$.  Notice that for $i_1 \neq 0$ 
\[
N=1+(p^2-p+1)^{i_1}+(p^2-p+1)^{2i_1}+\cdots+(p^2-p+1)^{(p-1)i_1}.
\]
We can write $p^2-p+1=wp+1$, where $w=p-1$. For any positive integer $k$, $$(1+wp)^k=\displaystyle \sum_{t=0}^{k}  \binom{k}{t} (wp)^t \equiv 1+kwp \pmod {p^2}.$$Thus, 
\[
N \equiv \displaystyle p+i_1pw\sum_{s=1}^{p-1}s \equiv p+i_1pw \frac{p(p-1)}{2} \equiv p \pmod {p^2}.
\]
 Hence,  $(b^{j_1}a^{i_1})^p=I$ if and only if $j_1N+i_1p \equiv 0 \pmod {p^2}$ if and only if $j_1+i_1 \equiv 0 \pmod p$. Consequently $|b^{j_1}a^{i_1}|=p^2$ implies $j_1+i_1 \not \equiv 0 \pmod  p$. Similarly, for $|b^{j_2}a^{i_2}|=p^2$, we can prove that $j_2+i_2 \not \equiv 0 \pmod  p$.

It remains to show that $i_2\equiv i_1-1 \pmod  p$ and $j_2 \equiv j_1+1 \pmod  p$. After raising $p$-th power to both sides of the equations in System \eqref{s41} and using $a^p=b^p$, we get $(b^{j_1}a^{i_1})^p=(b^{j_2}a^{i_2})^p$. This gives us  
\[
b^{j_1N+i_1p}=b^{j_2N'+i_2p},
\]
where $$
N'=\begin{cases} \dfrac{(p^2-p+1)^{i_2p}-1}{(p^2-p+1)^{i_2}-1} & \text{ if } i_2 \neq 0,\\p 
 & \text{ if } i_2=0.\end{cases}$$
  However, we know that equality $b^{j_1N+i_1p}=b^{j_2N'+i_2p}$ holds if and only if  $j_1N+i_1p \equiv j_2N'+i_2p \pmod {p^2}$.
 From the above discussion and definitions of $N$ and $N'$, it is clear that $N \equiv p \equiv N' \pmod {p^2}$ for $0 \leq i_1,i_2 \leq p-1$. Thus,  $j_1N+i_1p \equiv j_2N'+i_2p \pmod {p^2}$ if and only if $j_1+i_1 \equiv j_2+i_2 \pmod p$ or equivalently, $i_1-i_2 \equiv j_2-j_1 \pmod p$. Let $i_1-i_2 \equiv j_2-j_1 \equiv r \pmod p$ for some $0 \leq r \leq p-1$. It allows us to write $i_1=i_2+r+pt_1$ and $j_1=j_2-r+pt_2$ for some integers $t_1$ and $t_2$.
 
 Furthermore, we also have $ab=b^{p^2-p+1}a$ or equivalently $ab^{p+1}=ba$ (as $a^p=b^p$) and therefore using System \eqref{s41}, we obtain    
 \begin{equation}\label{e43}
 h_{a,b}:=b^{j_1}a^{i_1}(b^{j_2}a^{i_2})^{p+1}=b^{j_2}a^{i_2}b^{j_1}a^{i_1}=:g_{a,b}.
 \end{equation}

We now compute $h_{a,b}$.
%Using properties of $a$ and $b$ \textcolor{blue}{which properties??} and the fact that \textcolor{red}{$G^{p}:=\{a^p \mid a \in G\} \subseteq  \mathbb{Z}(G)$,} where $\mathbb{Z}(G)$ is the center of $G$, we compute
Notice that 
\[
h_{a,b} = b^{j_1}a^{i_1}(b^{j_2}a^{i_2})^{p+1} 
= (b^{j_2}a^{i_2})^{p} b^{j_1}a^{i_1}b^{j_2}a^{i_2},
\]
as $G^{p} := \{x^{p} \mid x \in G\} \subset \mathbb{Z}(G)$,
where $ \mathbb{Z}(G) $ denotes the center of $ G $, 
and the inclusion holds because $a^{p} = b^{p}$. Substituting the expressions for $ i_1$ and $ j_1 $ and using again the fact that $G^{p} \subset \mathbb{Z}(G) $, we obtain  
\[
h_{a,b}=(b^{j_2}a^{i_2})^{p} b^{j_1}a^{i_1}b^{j_2}a^{i_2}
= (b^{j_2}a^{i_2})^{p}b^{j_2-r}a^{i_2+r}b^{j_2}b^{pt_2}a^{i_2}a^{pt_1}.
\]

Further we employ the relation $ab=b^{p^2-p+1}a=b^{-p+1}a$ repeatedly to obtain $$a^rb^{j_2}=b^{j_2(-p+1)^{r}}a^r.$$Hence, 
\begin{equation*}
\begin{split}
h_{a,b}= (b^{j_2}a^{i_2})^{p}b^{j_2-r}a^{i_2+r}b^{j_2}b^{pt_2}a^{i_2}a^{pt_1}={(b^{j_2}a^{i_2})^{p}b^{j_2-r}a^{i_2}b^{j_2(-p+1)^{r}}b^{pt_2}a^{i_2+r+pt_1}}
%\\&={(b^{j_2}a^{i_2})^{p}b^{j_2-r}b^{j_2(-p+1)^r-j_2}a^{i_2}b^{j_2}b^{pt_2}a^{i_2+r+pt_1}}
\end{split}
\end{equation*}
as $b^{pt_2}\in \mathbb{Z}(G)$.

Moreover, observe that
\[
b^{j_2(-p+1)^r - j_2} = b^{\, j_2 \displaystyle\sum_{k=1}^{r}\binom{r}{k}(-p)^k} \in \mathbb{Z}(G),
\]
since $ p \Bigm| j_2 \left(\displaystyle\sum_{k=1}^{r}\binom{r}{k}(-p)^k\right)$ and $ G^{p} \subset \mathbb{Z}(G)$.
This allows us to rewrite
\[
\begin{aligned}
h_{a,b}
&= (b^{j_2}a^{i_2})^{p} b^{j_2-r} a^{i_2} b^{j_2(-p+1)^{r}} b^{-j_2}b^{j_2}b^{pt_2} a^{i_2+r+pt_1} \\
&= (b^{j_2}a^{i_2})^{p} b^{j_2-r} b^{j_2(-p+1)^r - j_2} a^{i_2} b^{j_2} b^{pt_2} a^{i_2+r+pt_1} \\
&= (b^{j_2}a^{i_2})^{p} b^{j_2-r} b^{j_2(-p+1)^r - j_2} a^{i_2} b^{r} b^{pt_2 + j_2 - r} a^{i_2+r+pt_1}.
\end{aligned}
\]
Since \( a^{i_2}b^{r} = b^{r(-p+1)^{i_2}}a^{i_2} \), therefore
\[
\begin{aligned}
h_{a,b}
&= (b^{j_2}a^{i_2})^{p} b^{j_2-r} b^{j_2(-p+1)^r - j_2} b^{r(-p+1)^{i_2}} a^{i_2} b^{pt_2 + j_2 - r} a^{i_2+r+pt_1} \\
&= (b^{j_2}a^{i_2})^{p} b^{j_2(-p+1)^r - j_2} b^{r(-p+1)^{i_2} - r} b^{j_2} a^{i_2} b^{j_1} a^{i_1}.
\end{aligned}
\]
As $$(b^{j_2}a^{i_2})^{p}=b^{j_2N'+i_2p},~ j_2(-p+1)^r-j_2 =j_2 \left(\sum_{k=1}^{r} \binom{r}{k}(-p)^k\right) \text{ and } r(-p+1)^{i_2}-r=r\left(\sum_{k=1}^{i_2}\binom{i_2}{k}(-p)^k\right),$$
consequently,
\[
h_{a,b}={(b^{j_2}a^{i_2})^{p}b^{j_2(-p+1)^r-j_2}b^{r(-p+1)^{i_2}-r}b^{j_2}a^{i_2}b^{j_1}a^{i_1}}=b^{u'}g_{a,b,}
\]
where 
\[
u'=j_2N'+i_2p+j_2\left(\displaystyle\sum_{k=1}^{r}\binom{r}{k}(-p)^k\right)+r\left(\displaystyle\sum_{k=1}^{i_2}\binom{i_2}{k}(-p)^k\right).
\]


In Equation \eqref{e43}, we have $h_{a,b}=g_{a,b}$, thus it follows that 
\[j_2N'+i_2p+j_2 \left(\sum_{k=1}^{r}\binom{r}{k}(-p)^k\right)+r\left(\sum_{k=1}^{i_2}\binom{i_2}{k}(-p)^k\right) \equiv j_2N'+i_2p-rp(j_2+i_2) \equiv 0 \pmod {p^2}.\] 
%\textcolor{blue}{How, thus came here?? at above you are writing we also have idnetity. here you are writnig thus, and how after thus N and N' came in picture and how follwing equalties you wrote???} 
Moreover, $$j_2N'+i_2p-rp(j_2+i_2) \equiv j_2p+i_2p-rp(j_2+i_2) \equiv 0 \pmod {p^2}$$ 
as $N' \equiv p \pmod {p^2}.$
It is straightforward to see that the congruence
\[
j_2p + i_2p - rp(j_2 + i_2) \equiv 0 \pmod{p^2}
\]
holds if and only if
$
(1 - r)(j_2 + i_2) \equiv 0 \pmod{p}.
$
Since \( (j_2 + i_2) \not\equiv 0 \pmod{p} \), it follows that \( 1 - r \equiv 0 \pmod{p} \), i.e., \( r = 1 \).
Hence, \( i_2 \equiv i_1 - 1 \pmod{p} \) and \( j_2 \equiv j_1 + 1 \pmod{p} \), as required.
 

Conversely, assume that \begin{equation*}
\begin{split}
& j_1+i_1 \not \equiv  0  \pmod p,\\& 
  j_2 +i_2\not \equiv  0 \pmod p,
\end{split}
\end{equation*}
 $i_2\equiv i_1-1 \pmod  p$ and $j_2 \equiv j_1+1 \pmod  p$. We have to show that $\mathcal{N}(i_1,j_1,i_2,j_2) \neq \emptyset$.  
For any $u \in \{0,1,\ldots,p-1\}$, $v \in \{0,1,\ldots,p^2-1\}$ and fixed $ 0 \leq t \leq p^3-1$, consider 
 \begin{equation}\label{tempe3}
 g(c_{v+up^2}):=\alpha^{v}\beta^{u}(c_t),
 \end{equation}
 where $\alpha:=b^{j_1}a^{i_1}$ and $\beta:=b^{j_2}a^{i_2}$. We shall show that $g$ is a permutation and $g \in \mathcal{N}(i_1,j_1,i_2,j_2)$. From the above discussion, we have $(b^{j_1}a^{i_1})^p=b^{j_1N+i_1p}$ and $ (b^{j_2}a^{i_2})^p=b^{j_2N'+i_2p}$, where 
$$
\begin{aligned}
N=&\begin{cases} \dfrac{(p^2-p+1)^{i_1p}-1}{(p^2-p+1)^{i_1}-1} & \text{ if } i_1 \neq 0,\\p 
 & \text{ if } i_1=0, \end{cases}\\
N'=&\begin{cases} \dfrac{(p^2-p+1)^{i_2p}-1}{(p^2-p+1)^{i_2}-1} & \text{ if } i_2 \neq 0,\\p 
 & \text{ if } i_2=0\end{cases}
 \end{aligned}$$
  and $N \equiv p \equiv N' \pmod {p^2}$. We first show that $|\alpha|=p^2=|\beta|$.
  Suppose, for contradiction, that $|\alpha|\neq p^{2}$. Then $\alpha^{p}=I$ as $|\alpha| \Bigm| p^2$, and hence
\[
(b^{j_1}a^{i_1})^{p}=b^{\,j_1N+i_1p}=I.
\]
This gives $j_1N+i_1p \equiv 0 \pmod{p^{2}}$, and since $N \equiv p \pmod{p^{2}}$,
this is equivalent to
$
j_1p+i_1p \equiv 0 \pmod{p^{2}}.
$
The congruence $j_1p+i_1p \equiv 0 \pmod{p^{2}}$ then yields
\[
j_1+i_1 \equiv 0 \pmod{p},
\]
which is a contradiction. Similarly, we can show $|\beta|=p^2$. Now we prove that $\alpha^p=\beta^p$. Since $i_2\equiv i_1-1 \pmod  p$ and $j_2 \equiv j_1+1 \pmod  p$, we have $i_2+j_2 \equiv i_1+j_1 \pmod  p$, which gives  $i_2p+j_2p \equiv i_1p+j_1p \pmod {p^2}$. Using $N \equiv p \equiv N' \pmod {p^2}$, it follows that  $j_1N+i_1p \equiv j_2N'+i_2p \pmod {p^2}$. Therefore, 
$$\alpha^p=(b^{j_1}a^{i_1})^p=b^{j_1N+i_1p}=b^{j_2N'+i_2p}=(b^{j_2}a^{i_2})^p=\beta^p.$$
Next, we prove that $\alpha \beta=\beta^{p^2-p+1}\alpha$. Since $i_2\equiv i_1-1 \pmod  p$ and $j_2 \equiv j_1+1 \pmod  p$, we may write $i_2=i_1-1+pt_1$ and $j_2=j_1+1+pt_2$ for some integers $t_1,t_2$. Now,
  \[
  \begin{aligned}
  \alpha \beta=&b^{j_1}a^{i_1}b^{j_1+1+pt_2}a^{i_1-1+pt_1}
  \\=&b^{j_1+pt_2}a^{i_1}ba^{pt_1}b^{j_1}a^{i_1-1}
    \end{aligned}
  \]
 as $b^{pt_2},a^{pt_1} \in \mathbb{Z}(G)$. Using the relation $ab=b^{p^2-p+1}a=b^{-p+1}a$, we obtain $a^{i_1}b=b^{(-p+1)^{i_1}}a^{i_1}$. Hence, 
 \[
 \alpha \beta=b^{j_1+pt_2}b^{(-p+1)^{i_1}}a^{i_1}a^{pt_1}b^{j_1}a^{i_1-1}=b^{j_1+pt_2}b^{(-p+1)^{i_1}}a^{i_1+pt_1-1}ab^{j_1}a^{i_1-1}.
 \]
 Again, from $ab=b^{-p+1}a$, we have $ab^{j_1}=b^{j_1(-p+1)}a$, which implies 
 \begin{equation*}
     \alpha \beta=b^{j_1+pt_2}b^{(-p+1)^{i_1}}a^{i_1+pt_1-1}b^{j_1(-p+1)}a^{i_1}.
 \end{equation*}
 Using $(-p+1)^{i_1} \equiv -pi_1+1 \pmod {p^2}$, $|b|=p^2$ and $b^{-j_1p} \in \mathbb{Z}(G)$,
\begin{equation}\label{tempe2}
     \alpha \beta=b^{-p(i_1+j_1)}b^{j_1+1+pt_2}a^{i_1+pt_1-1}b^{j_1}a^{i_1}=b^{-p(i_1+j_1)}\beta \alpha.
 \end{equation}
 The congruence  $i_1+j_1 \equiv i_2+j_2 \pmod  p$ implies $p(i_1+j_1) \equiv p(i_2+j_2) \pmod {p^2}$ and therefore $$b^{-p(i_1+j_1)}=b^{-p(i_2+j_2)}=(b^{j_2N'+pi_2})^{-1}=(b^{j_2}a^{i_2})^{-p}=\beta^{-p}=\beta^{p^2-p}$$ as $N'\equiv p \pmod {p^2}$. Substituting this into \eqref{tempe2} gives $$\alpha \beta=\beta^{p^2-p+1}\alpha.$$ Thus, $G=\{\alpha^{v}\beta^{u} \mid 0 \leq u \leq p-1, 0 \leq v \leq p^2-1\}$. In Theorem~\ref{T31}, we established that every non-identity permutation in $G$
has no fixed point in $\mathbb{F}_q$. Consequently, no two distinct
permutations in $G$ can take the same value at any element of $\mathbb{F}_q$,
as $G$ is a subgroup of $\mathfrak{S}_q$. Therefore,
\[
\{\alpha^{v}\beta^{u}(c_t)\mid 0 \le u \le p-1,\; 0 \le v \le p^{2}-1\}
 = \mathbb{F}_q,
\]
which shows that $g$ is indeed a permutation.

 
 
 Now we will show that $g \in \mathcal{N}(i_1,j_1,i_2,j_2)$, that is, 
  \begin{equation}\label{s42}
  \begin{split}
  gag^{-1}=b^{j_1}a^{i_1}=\alpha \\
  gbg^{-1}=b^{j_2}a^{i_2}=\beta.
  \end{split}
  \end{equation} 
  Since $g$ is a permutation and $\{0,1,\ldots,p^3-1\}=\{v+up^2 \mid 0 \leq u \leq p-1, 0 \leq v \leq p^2-1\}$, it suffices to verify both equations in System \eqref{s42} for all the elements of the form $g(c_{v+up^2})$, where $u \in \{0,1,\ldots,p-1\}$ and $v \in \{0,1,\ldots,p^2-1\}$.  From Equation \eqref{e21}, we have $a(c_{v+up^2})=c_{up^2+(v+1) \pmod {p^2}}$ thus, 
  $$gag^{-1}(g(c_{v+up^2}))=ga(c_{v+up^2})=g(c_{up^2+(v+1) \pmod {p^2}}).$$ 
Further, Equation \eqref{tempe3} implies that $g(c_{v+up^2})=\alpha^{v}\beta^{u}(c_t)$ and $$g(c_{(v+1) \pmod {p^2}+up^2})=\alpha^{(v+1) \pmod {p^2}}\beta^{u}(c_t)=\alpha^{v+1}\beta^{u}(c_t)$$ as $|\alpha|=p^2$.
Thus,
 \[
 \alpha(g(c_{v+up^2}))= \alpha(\alpha^{v}\beta^{u}(c_t))=\alpha^{v+1}\beta^{u}(c_t)=g(c_{(v+1) \pmod {p^2}+up^2})=gag^{-1}(g(c_{v+up^2})).
 \]
 Hence, $g$ satisfies the first equation of the System \eqref{s42}. 
 
 For the second equation of the System \eqref{s42}, assume that $v \equiv k \pmod  p$ for some $k \in \{0,1,\ldots,p-1\}$, which immediately yields  $v+up^2 \equiv k \pmod  p$ and thus $b(c_{v+up^2})=b_k^{1+kp}(c_{v+up^2})$ as from the definition of $b$, $b=\displaystyle \prod_{\ell=0}^{p-1}b_{\ell}^{1+\ell p}$ and Equation \eqref{e22} implies $b_{\ell}(c_t)=c_t$ whenever $t \not \equiv \ell \pmod  p$. %Therefore, 
 %\[
%  gbg^{-1}(g(c_{v+up^2}))= g(b(c_{v+up^2}))=g(b_k^{1+kp}(c_{v+up^2})).
 %\]
 Since $a^p=b^p$, we get $a^{kp}(c_{v+up^2})=b^{kp}(c_{v+up^2})=b_k^{kp(1+kp)}(c_{v+up^2})=b_k^{kp}(c_{v+up^2}).$  Furthermore, using Equation \eqref{tempe4}, 
 $$a^{kp}(c_{v+up^2})=c_{up^2+(v+kp)\pmod {p^2}}$$
as $up^2\leq v+up^2 \leq p^2-1+up^2 =(u+1)p^2-1$.

Thus, 
\begin{equation}\label{R43}
b(c_{v+up^2})=b_k^{1+kp}(c_{v+up^2})=b_k(a^{kp}(c_{v+up^2}))=b_k(c_{up^2+(v+kp)\pmod {p^2}}).
\end{equation}
Now we verify the second equation of System \eqref{s42} into two cases depending on the value of $u$.
 
 \textbf{Case 1:} In this case, we consider $u<p-1$. Note that $up^2+(v+kp)\pmod {p^2} \leq (p-2)p^2+p^2-1=p^3-p^2-1<p^3-p^2.$ Therefore using Equation \eqref{e22} in Equation \eqref{R43}, we have 
 \[
 b(c_{v+up^2})=b_k^{1+kp}(c_{v+up^2})=b_{k}(c_{up^2+(v+kp)\pmod {p^2}})=c_{(u+1)p^2+(v+kp)\pmod {p^2}}.
 \]
 This gives 
  \[
  gbg^{-1}(g(c_{v+up^2})) = g(b(c_{v+up^2}))=g(c_{(u+1)p^2+(v+kp)\pmod {p^2}}).
 \]
Moreover, from the definition of $g$ in Equation \eqref{tempe3}, we obtain
\[
g(c_{v+up^2})=\alpha^{v}\beta^{u}(c_t) \text{ and } g(c_{(u+1)p^2+(v+kp)\pmod {p^2}})=\alpha^{(v+kp)\pmod {p^2}}\beta^{u+1}(c_t)=\alpha^{v+kp}\beta^{u+1}(c_t), 
\]
as $|\alpha|=p^2$.
Using $\alpha \beta=\beta^{p^2-p+1}\alpha$, $|\beta|=p^2$ and $\alpha^p=\beta^p$, we can easily get $\beta \alpha=\alpha \beta^{p+1}$. The equality $\beta \alpha=\alpha \beta^{p+1}$ gives $\beta \alpha^v=\alpha^v\beta^{(p+1)^v}=\alpha^v\beta^{vp+1}$ as $(p+1)^v \equiv 1+vp \pmod {p^2}$ and $|\beta|=p^2$. Therefore,
\begin{equation*}
\beta(g(c_{v+up^2}))=\beta(\alpha^{v}\beta^{u}(c_t)) = \alpha^v \beta^{vp+1}\beta^{u}(c_t)=\alpha^{v+vp}\beta^{u+1}(c_t)=\alpha^{v+kp}\beta^{u+1}(c_t)
\end{equation*}
as $\alpha^p=\beta^p$ and $vp \equiv kp \pmod {p^2}$. Thus
 \[ gbg^{-1}(g(c_{v+up^2}))=g(c_{(u+1)p^2+(v+kp)\pmod {p^2}})=\alpha^{v+kp}\beta^{u+1}(c_t)=\beta(g(c_{v+up^2})).\]
 %Thus, $gbg^{-1}(g(c_{v+up^2}))=\beta(g(c_{v+up^2}))$.

\textbf{Case 2:} When $u=p-1$. We divide this case into following two subcases.

\textbf{Subcase 2.1:} In this subcase, we consider $0 \leq (v+kp)\pmod {p^2}<p^2-p$. So, we have $p^3-p^2 \leq up^2+(v+kp)\pmod {p^2} <p^3-p$. By applying Equation \eqref{e22} to \eqref{R43}, we have 
\[
b(c_{v+up^2})=b_k^{1+kp}(c_{v+up^2})=b_k(c_{up^2+(v+kp)\pmod {p^2}})=c_{p+(v+kp)\pmod {p^2}}.
\]
Thus, 
  \[
  gbg^{-1}(g(c_{v+up^2}))= g(b(c_{v+up^2}))=g(c_{p+(v+kp)\pmod {p^2}}).
 \]
From the definition of $g$,
\[
gbg^{-1}(g(c_{v+up^2}))=g(c_{p+(v+kp)\pmod {p^2}})=\alpha^{p+v+kp}(c_t), 
\]
as $0<p+(v+kp)\pmod {p^2}<p+p^2-p= p^2$, that is, $0<p+(v+kp)\pmod {p^2} \leq p^2-1$.
On the right side of the second equation in System \eqref{s42},
\[
\beta(g(c_{v+up^2}))=\beta(\alpha^{v}\beta^{u}(c_t))= \alpha^v \beta^{vp+1} \beta^{u}(c_t)=\alpha^{v+vp+p}(c_t)=\alpha^{v+kp+p}(c_t)
\]
as $g(c_{v+up^2})=\alpha^{v}\beta^{u}(c_t)$, $vp \equiv kp \pmod {p^2}$, $\alpha^p=\beta^p$ and $u+1=p$. Therefore, $gbg^{-1}(g(c_{v+up^2}))=\beta(g(c_{v+up^2}))$.

\textbf{Subcase 2.2:} In this subcase, we assume that $p^2-p \leq (v+kp)\pmod {p^2} <p^2$. This subcase can be processed similarly to Subcase 2.1. 

From the above cases, we conclude that $g \in \mathcal{N}(i_1,j_1,i_2,j_2)$ and so $\mathcal{N}(i_1,j_1,i_2,j_2) \neq \emptyset$. Observe that, for fixed $\alpha = b^{j_1}a^{i_1}$ and $\beta = b^{j_2}a^{i_2}$, the function $g$ is completely determined by the choice of $c_t$, and there are $p^{3}$ possible choices for $c_t$. Moreover, from Equation \eqref{tempe3}, we have $g(c_0)=c_t$, which ensures that distinct choices of $c_t$ produce distinct functions $g$. Hence,
$|\mathcal{N}(i_1,j_1,i_2,j_2)| \ge p^{3}.$
Finally, using Lemma \ref{Atmost_lemma_T31}, we have $|\mathcal{N}(i_1,j_1,i_2,j_2)|=p^3$.
\end{proof}

\end{document}
